\section{Method}
\label{sec:method}

\subsection{Problem Formulation and Design Overview}
\label{sec:method_overview}

We study discrete action tokenization for bimanual manipulation.  An action
chunk is denoted by
\begin{equation}
    \mathbf{A}
    = [\mathbf{A}^{L},\mathbf{A}^{R}]
    \in \mathbb{R}^{H\times 2d_a},
\end{equation}
where $H$ is the action horizon, $d_a$ is the action dimension of one arm, and
$L$ and $R$ indicate the left and right arms, respectively.  Our goal is to
learn a tokenizer $\mathcal{T}$ and a detokenizer $\mathcal{D}$ such that
\begin{equation}
    (\mathbf{y}^{L},\mathbf{y}^{R}) = \mathcal{T}(\mathbf{A}),
    \qquad
    \widehat{\mathbf{A}}
    = \mathcal{D}(\mathbf{y}^{L},\mathbf{y}^{R}),
    \label{eq:btt_mapping}
\end{equation}
where each arm is represented by $M\ll H$ categorical tokens,
$\mathbf{y}^{r}\in\{0,\ldots,K-1\}^{M}$ for
$r\in\{L,R\}$.  In our default configuration,
$(H,d_a,M,K)=(32,7,8,1920)$, resulting in 16 tokens for a 32-step,
14-dimensional bimanual action chunk.

A bimanual tokenizer must preserve two complementary structures.  Each arm
obeys similar local kinematics, favoring arm-shared representations, while task
completion depends on precisely timed inter-arm coordination, favoring joint
modeling.  Treating the full 14-dimensional trajectory as a monolithic signal
obscures the former structure; tokenizing the two arms independently discards
the latter.  We therefore propose \emph{Bimanual Temporal Tokenization (BTT)},
which factorizes token identity by arm but couples the two streams on both
sides of the discrete bottleneck.  Specifically, BTT consists of: (i) a shared
causal temporal encoder that extracts single-arm motion primitives;
(ii) symmetric causal cross-attention that injects coordination before
quantization; (iii) a shared finite scalar quantizer that produces separate
left- and right-arm token streams; and (iv) a query-based decoder that pairs
the two dequantized streams at corresponding latent timesteps and jointly
reconstructs both arms.  This design keeps arm tokens directly addressable
without forcing the decoder to infer coordination from two unrelated latent
sequences.

\subsection{Canonical Bimanual Action Space}
\label{sec:btt_canonicalization}

Physical left and right arms commonly use mirrored joint conventions.  Naively
sharing an encoder across their raw coordinates would therefore assign
different meanings to the same feature channel.  Before encoding, we map both
arms into a common kinematic convention.  Let
$\mathbf{s}\in\{-1,+1\}^{d_a}$ be a fixed sign vector.  For timestep $t$, we
define
\begin{equation}
    \overline{\mathbf{a}}^{L}_{t}=\mathbf{a}^{L}_{t},
    \qquad
    \overline{\mathbf{a}}^{R}_{t}
    =\mathbf{s}\odot\mathbf{a}^{R}_{t},
    \label{eq:btt_canonicalization}
\end{equation}
where $\odot$ denotes element-wise multiplication.  Because
$s_j\in\{-1,+1\}$, this transformation is self-inverse.

We then normalize both canonical arms with the \emph{same} joint-wise bounds.
The bounds $\boldsymbol{\ell},\mathbf{u}\in\mathbb{R}^{d_a}$ are computed by
pooling the left-arm data and canonicalized right-arm data.  The normalized
action is
\begin{equation}
    \mathbf{x}^{r}_{t}
    = 2\frac{\overline{\mathbf{a}}^{r}_{t}-\boldsymbol{\ell}}
             {\mathbf{u}-\boldsymbol{\ell}}-1,
    \qquad r\in\{L,R\},
    \label{eq:btt_normalization}
\end{equation}
with all operations applied element-wise.  Sharing both the coordinate
convention and the normalization statistics makes weight sharing structural:
a motion executed by either arm is presented to the encoder in the same
feature space.

\subsection{Symmetry-Aware Causal Encoder}
\label{sec:btt_encoder}

The BTT encoder first models per-arm dynamics with shared parameters and then
models cross-arm coordination.  This separation provides an arm-swap
equivariance in canonical coordinates while allowing every token to carry
information about the other arm.

\paragraph{Shared causal temporal compression.}
For each arm $r$, a shared linear layer lifts
$\mathbf{X}^{r}_{0}=\mathbf{x}^{r}_{1:H}$ from $d_a$ to the model dimension
$d$.  We subsequently apply two weight-shared causal convolutional stages,
\begin{equation}
    \mathbf{X}^{r}_{\ell}
    = \operatorname{Mish}\!\left(
        \operatorname{LN}\!\left(
        \operatorname{CConv}_{k_{\ell},s_{\ell}}
        (\mathbf{X}^{r}_{\ell-1})\right)\right),
    \quad \ell\in\{1,2\}.
    \label{eq:btt_causal_conv}
\end{equation}
The convolution uses left-only padding $k_{\ell}-s_{\ell}$.  Consequently,
output position $i$ is aligned with the end of its stride interval and cannot
depend on a later input.  With kernels $(5,3)$ and strides $(2,2)$, the two
stages reduce $H=32$ steps to $M=8$ latent positions.  Layer normalization is
applied independently at each timestep, avoiding temporal-statistics leakage.

We add a fixed sinusoidal position embedding $\mathbf{P}^{M}$ after temporal
compression and use a shared pre-normalized, causally masked self-attention
layer:
\begin{equation}
    \mathbf{S}^{r}
    = \operatorname{CSA}_{\theta}
      (\mathbf{X}^{r}_{2}+\mathbf{P}^{M}),
    \qquad r\in\{L,R\}.
    \label{eq:btt_single_arm}
\end{equation}
The same projection, convolutions, normalization layers, and self-attention
parameters $\theta$ are used for both arms.  Thus, BTT learns one reusable
model of single-arm dynamics instead of duplicating arm-specific encoders.

\paragraph{Symmetric causal coordination.}
Independent arm features are insufficient for motions such as handovers,
cooperative grasping, or synchronized placement.  We therefore exchange
information between the two streams before discretization.  Using one shared
cross-attention operator in both directions, we compute
\begin{align}
    \mathbf{C}^{L}
    &= \mathbf{S}^{L}
       + \operatorname{CA}_{\phi}
       \!\left(\operatorname{LN}(\mathbf{S}^{L}),
               \operatorname{LN}(\mathbf{S}^{R}),
               \operatorname{LN}(\mathbf{S}^{R});\mathbf{M}_{\mathrm{c}}\right),
       \nonumber\\
    \mathbf{C}^{R}
    &= \mathbf{S}^{R}
       + \operatorname{CA}_{\phi}
       \!\left(\operatorname{LN}(\mathbf{S}^{R}),
               \operatorname{LN}(\mathbf{S}^{L}),
               \operatorname{LN}(\mathbf{S}^{L});\mathbf{M}_{\mathrm{c}}\right),
    \label{eq:btt_cross_attention}
\end{align}
where $\mathbf{M}_{\mathrm{c}}[i,j]=0$ for $j\leq i$ and $-\infty$
otherwise.  Both updates are evaluated in parallel from the original
$(\mathbf{S}^{L},\mathbf{S}^{R})$, and the attention parameters $\phi$ are
shared across directions.  Hence swapping the two canonical input arms swaps
the two encoder outputs.  The causal mask further ensures that coordinated
latent position $i$ only summarizes motion observed up to its aligned action
timestep, rather than leaking future actions into earlier tokens.

\subsection{Arm-Factorized Finite Scalar Quantization}
\label{sec:btt_quantization}

The coordinated features remain separated by arm.  A shared projection maps
each feature vector to a compact latent,
\begin{equation}
    \mathbf{v}^{r}_{i}
    = \mathbf{W}_{q}\mathbf{C}^{r}_{i}+\mathbf{b}_{q}
    \in\mathbb{R}^{d_z},
    \qquad i=1,\ldots,M.
    \label{eq:btt_prequant}
\end{equation}
We discretize every scalar channel with finite scalar quantization (FSQ).  For
per-channel level counts $\mathbf{L}=(L_1,\ldots,L_{d_z})$, FSQ smoothly bounds
each coordinate, rounds it to the corresponding finite grid, and uses a
straight-through estimator for the rounding operation.  If
$n^{r}_{i,j}\in\{0,\ldots,L_j-1\}$ denotes the resulting digit, its mixed-radix
token ID is
\begin{equation}
    y^{r}_{i}
    = \sum_{j=1}^{d_z} n^{r}_{i,j}
      \prod_{m<j}L_m,
    \qquad
    K=\prod_{j=1}^{d_z}L_j.
    \label{eq:btt_fsq}
\end{equation}
We use $d_z=4$ and levels $(8,8,6,5)$, giving $K=1920$ possible symbols.  The
pre-quantization projection and FSQ grid are shared across arms, so a token ID
has the same quantization semantics for the left and right streams.

Crucially, BTT does not collapse an aligned left--right pair into a single
categorical variable.  It emits two length-$M$ sequences and serializes them as
\begin{equation}
    \mathbf{y}
    = [y^{L}_{1},\ldots,y^{L}_{M},
       y^{R}_{1},\ldots,y^{R}_{M}].
    \label{eq:btt_serialization}
\end{equation}
This retains an ordinary $K$-way vocabulary and exposes the arm structure to a
downstream policy, while the pre-quantization cross-attention ensures that the
symbols are not learned in isolation.

\subsection{Timestep-Paired Joint Decoder}
\label{sec:btt_decoder}

The decoder deliberately complements the factorized token interface.  Given
token IDs, we first recover their FSQ embeddings
$\widehat{\mathbf{z}}^{L},\widehat{\mathbf{z}}^{R}
\in\mathbb{R}^{M\times d_z}$.  Rather than concatenating the two arm streams
along time or decoding each arm separately, we pair embeddings at every
aligned latent timestep:
\begin{equation}
    \mathbf{m}_{i}
    = \mathbf{W}_{m}
      [\widehat{\mathbf{z}}^{L}_{i}\,\Vert\,
       \widehat{\mathbf{z}}^{R}_{i}]
      +\mathbf{b}_{m}
    \in\mathbb{R}^{d},
    \quad i=1,\ldots,M,
    \label{eq:btt_paired_memory}
\end{equation}
where $\Vert$ denotes feature concatenation.  The resulting $M$ memory slots
explicitly preserve which two arm codes describe the same temporal segment.
Pairing is performed after dequantization; it therefore captures the joint
state without requiring a $K^2$-sized pair vocabulary.

To reconstruct an $H$-step action chunk, we instantiate $H$ zero-valued queries
and add fixed sinusoidal position embeddings
$\mathbf{Q}_{0}=\mathbf{P}^{H}$.  A Transformer decoder maps the queries and
paired memory to action features,
\begin{equation}
    \mathbf{F}
    = \operatorname{TrDec}_{\psi}
      (\mathbf{Q}_{0},\mathbf{m}_{1:M};\mathbf{M}_{\mathrm{c}}),
    \qquad
    \widehat{\mathbf{X}}
    = \mathbf{F}\mathbf{W}_{o}^{\top}+\mathbf{b}_{o},
    \label{eq:btt_decode}
\end{equation}
where $\widehat{\mathbf{X}}\in\mathbb{R}^{H\times 2d_a}$.  Self-attention among
the action queries is causally masked, whereas cross-attention is unmasked so
that every output timestep can access all $M$ paired memory slots.  No previous
ground-truth action is supplied to the decoder, and all $H$ outputs are
reconstructed in one forward pass.  A single output head predicts the complete
14-dimensional bimanual vector, forcing the decoder to resolve coordination
jointly rather than with two independent reconstruction heads.

Finally, we split $\widehat{\mathbf{X}}$ into two normalized arm trajectories,
invert Eq.~\eqref{eq:btt_normalization}, and apply the same sign transform in
Eq.~\eqref{eq:btt_canonicalization} to the reconstructed right arm.  This
returns $\widehat{\mathbf{A}}$ in the original physical action convention.

\subsection{Learning Objective and Policy Interface}
\label{sec:btt_training}

We train the tokenizer end-to-end using reconstruction in the shared normalized
space.  For the set $\mathcal{V}$ of action chunks containing no padded
timesteps, the objective is
\begin{equation}
    \mathcal{L}_{\mathrm{rec}}
    = \frac{1}{|\mathcal{V}|H(2d_a)}
      \sum_{b\in\mathcal{V}}
      \left\|
      \widehat{\mathbf{X}}_{b}-\mathbf{X}_{b}
      \right\|_{F}^{2}.
    \label{eq:btt_objective}
\end{equation}
FSQ is parameter-free and is optimized through the straight-through estimator,
so BTT requires neither a learned codebook nor auxiliary commitment and
codebook losses.  The objective jointly trains the shared temporal encoder,
inter-arm attention, pre-quantization projection, paired decoder, and action
head.

After tokenizer pretraining, we freeze $\mathcal{T}$ and $\mathcal{D}$ and use
the 16 symbols in Eq.~\eqref{eq:btt_serialization} as action targets for a
vision-language-action policy.  The serialization naturally induces two
prediction blocks of sizes $(M,M)$: the policy predicts the left-arm block in
parallel from multimodal context $\mathbf{o}$ and then predicts the right-arm
block conditioned on the completed left block,
\begin{equation}
    p(\mathbf{y}\mid\mathbf{o})
    = \prod_{i=1}^{M}p(y^{L}_{i}\mid\mathbf{o})
      \prod_{i=1}^{M}p(y^{R}_{i}\mid\mathbf{o},\mathbf{y}^{L}).
    \label{eq:btt_policy_factorization}
\end{equation}
We optimize categorical cross-entropy with equal weight for the two blocks.
At inference, two blockwise decoding iterations produce the complete token
sequence, which the frozen BTT decoder converts into an executable bimanual
action chunk.  This interface preserves the computational advantage of
parallel prediction within each arm while exposing an explicit causal route
from one arm's plan to the other.

\paragraph{Architecture details.}
Unless stated otherwise, BTT uses model width $d=256$, eight attention heads,
and dropout $0.1$.  The query decoder contains six pre-normalized Transformer
layers with GELU feed-forward networks of width $4d$.  Both encoder and decoder
use fixed sinusoidal temporal embeddings.  For our 7-DoF arms, the right-arm
canonicalization vector is
$\mathbf{s}=(-1,1,1,1,-1,-1,1)$.
